\documentclass[10pt,a4paper]{amsart}
\usepackage[margin=19mm]{geometry}
\usepackage{amsmath,amssymb,mathtools}
\usepackage[scr=rsfs,cal=euler]{mathalfa}
\usepackage{xfrac}
\usepackage[unicode,hidelinks,bookmarks=false]{hyperref}
\newcommand{\ie}{i.e.\ }
\newcommand{\cf}{cf.\ }
\newcommand{\resp}{respectively\ }
\newcommand{\Subsection}{\subsection}
\providecommand{\nobreakdash}{\nobreak\hbox{-}}
\setlength{\emergencystretch}{2em}
\multlinegap0pt
\usepackage{bm}
\usepackage{bbm}
\usepackage{dsfont}
\usepackage{xspace}
\usepackage{cancel}
\usepackage{mathtools}
\usepackage{cleveref}
\usepackage{amsthm}
\makeatletter
\@ifundefined{theorem}{\newtheorem{theorem}{Theorem}}{}
\@ifundefined{lemma}{\newtheorem{lemma}[theorem]{Lemma}}{}
\makeatother
\makeatletter
\expandafter\ifx\csname ea\endcsname\relax
\newenvironment{ea}{
	\begin{equation}
	\begin{aligned}
}{
	\end{aligned}
	\end{equation}
	\ignorespacesafterend
}
\fi
\makeatother
\title[Diffusion-Reaction Epidemic Model with a Free Boundary]{Diffusion-Reaction Epidemic Model with a Free Boundary}
\author{Aesol Jeon, Ki-Ahm Lee}
\date{Source: arXiv:2511.07291v1 (CC BY 4.0)}
\begin{document}
\maketitle

\noindent\emph{Source and licence.} Diffusion-Reaction Epidemic Model with a Free Boundary. Version arXiv:2511.07291v1 (CC BY 4.0), distributed under the Creative Commons Attribution 4.0 International licence, as recorded in the arXiv licence metadata for this exact version and shown on the abstract page https://arxiv.org/abs/2511.07291v1. Full rights notice: licenses/arxiv-2511.07291-CC-BY-4.0.txt.\par\smallskip

\noindent\textit{Editorial scope.} Counted unit: the Lemma whose source label is \texttt{equation changing} in the article (source0.tex lines 1382--1409, source label \texttt{equation changing}), delivered together with the article's own complete original proof of it (source0.tex lines 1411--1528, one \texttt{proof} environment of 118 lines). No mathematics is written, repaired or re-derived: statement and proof are the article's own text line for line. Required context retained uncounted: free boundary SEIS (lines 342--351); initial cond (lines 353--358); straightening lem2 (lines 1352--1354). Every definition, display and label of those lines is retained unchanged. Omitted from this package and not redistributed: 1--341, 352--352, 359--1351, 1355--1381, 1410--1410, 1529--2114. Nothing inside a retained excerpt is omitted or altered: every retained body is delivered exactly as the article's lines are written, so no omission receipt of any kind accompanies this package. Citations: the retained excerpts carry no citation command at all, so no bibliography entry of the article is transcribed and no reference is redistributed. Reference adaptation: none was needed and none was made; every reference inside the delivered unit resolves inside it, so no unresolved reference or citation is produced. Printed slips kept literal: no printed word, symbol, hypothesis or number of the counted statement or of its proof was altered. Layout and class: the article's own class is \texttt{amsart} with packages including mathrsfs, mathtools, a4wide, graphicx, amssymb, epstopdf, enumerate, titletoc, hyperref, cleveref, esint, verbatim, enumitem, kotex; the wrapper is the lane's pinned \texttt{amsart} editorial wrapper at 10pt on A4, which restates the theorem-like environments the retained text needs and adds the article's own macro definitions that the retained text needs (none), with extra wrapper packages (bm, bbm, dsfont, xspace, cancel, mathtools, cleveref). The delivered numbering is the unit's own. Assets: no retained span contains a figure environment, a graphics-inclusion command, a PDF-page inclusion, a table or a TikZ picture; no image file is copied into this package, the package declares no asset dependency and no image file is redistributed with it. Licence: version arXiv:2511.07291v1 of this article is distributed under the Creative Commons Attribution 4.0 International licence, as recorded in the arXiv licence metadata for this exact version and shown on the abstract page https://arxiv.org/abs/2511.07291v1. A full rights notice is delivered at licenses/arxiv-2511.07291-CC-BY-4.0.txt. Characters: every retained excerpt is pure ASCII, so the delivered engine, XeLaTeX with its default Latin Modern text font, reports no missing character.
 	\begin{equation} \label{free boundary SEIS}
 	\begin{cases}
			S_t-d_1\Delta S=A-\alpha IS-\mu_1 S+r_2 I+r_1E, &(r,t)\in\Omega_\infty,\\
			E_t-d_2\Delta E=(1-p)\alpha IS-(\beta_1 +r_1+\mu_2) E, &(r,t) \in \Omega_h,\\
			I_t-d_3\Delta I=p\alpha IS+\beta_1 E-(r_2+\mu_3) I, &(r,t) \in \Omega_h,\\
			h'(t)=-\beta E_r(h(t),t)-\mu I_r(h(t),t), &t>0\\
			E(r,t)=I(r,t)=0, &(r,t)\notin \Omega_h,\\
			S_r(0,t)=E_r(0,t)=I_r(0,t)=0 &t>0,
	\end{cases}
	\end{equation}

	\begin{equation} \label{initial cond}
		\left \{ \begin{aligned}
			&h(0)=h_0>0, \\
			&S(r,0)=S_0,\ E(r,0)=E_0,\ I(r,0)=I_0, \qquad r\geq0, \\
		\end{aligned} \right.
	\end{equation}

\begin{lemma} \label{straightening lem2}
	Suppose that $\displaystyle \lim _{t \to \infty}h(t)=h_\infty$ for some $h_\infty >0$. The free boundary $h(t)$ of the system of equations \cref{free boundary SEIS} with \cref{initial cond} is diffeomorphic to $h_\infty$ for $d_1=d_2=d_3=d$.
\end{lemma}

\begin{lemma} \label{equation changing}
Suppose that $\displaystyle \lim _{t \to \infty}h(t)=h_\infty$ for some $h_\infty >0$. If $d_1=d_2=d_3=d$ for some $d>0$, the solution of the equations \cref{free boundary SEIS} of the Exposed and the Infected individuals are equal to the solution of the equations,
\begin{equation} \label{modified free boundary SEIS4}
	\left \{ \begin{aligned}
		u_t-Xd\Delta_s u-(Yd+Zh'(t)+Wd)u_s=(1-p)\alpha vw-(\beta_1+r_1+\mu_2)u,\\
		v_t-Xd\Delta_s u-(Yd+Zh'(t)+Wd)u_s=p\alpha vw+\beta_1 u-(r_2+\mu_3)v,
	\end{aligned} \right.
\end{equation}
for $0<s<h_\infty$ and $t>0$, and
\begin{equation} \label{modified free boundary SEIS4 cond}
	\left \{ \begin{aligned}
		&h'(t)=-\beta v_s(h_0,t)-\mu w_s(h_0,t), \quad &&s=h_\infty,\\
		&u(s,t)=v(s,t)=0, &&s\geq h_\infty,
	\end{aligned} \right.
\end{equation}
where
\begin{equation*}
	\begin{aligned}
		&\sqrt{X(h(t),s)}\coloneqq \frac{\partial s}{\partial r}=\frac{1}{1+\xi'(s)(h(t)-h_0)},
			\\
		&Y(h(t),s)\coloneqq \frac{\partial^2s}{\partial r^2}=-\frac{\xi''(s)(h(t)-h_0)}{[1+\xi'(s)(h(t)-h_0)]^3},
			\\
		&Z(h(t),s)\coloneqq -\frac{1}{h'(t)}\frac{\partial s}{\partial t}=\frac{\xi(s)}{1+\xi'(s)(h(t)-h_0)},
			\\
		&W(h(t),s)\coloneqq(n-1)X\frac{(s\xi'-\xi)(h(t)-h_0)}{s(s+\xi(s)(h(t)-h_0))}.\,
	\end{aligned}
\end{equation*}
\end{lemma}

\begin{proof}
By \cref{straightening lem2}, free boundary $h(t)$ is changed to $h_\infty$. Now, set 
\begin{equation*}
	\begin{aligned}
	E(r,t) = E(s+\xi(s)(h(t)-h_\infty),t)=:u(s,t),\\
	I(r,t) = I(s+\xi(s)(h(t)-h_\infty),t)=:v(s,t).
	\end{aligned}
\end{equation*}
If we differentiate both sides with respect to $s$,
\begin{equation*}
	\begin{aligned}
		(1+\xi'(s)(h(t)-h_\infty))E_r=u_s,\\
		(1+\xi'(s)(h(t)-h_\infty))I_r=v_s,
	\end{aligned}
\end{equation*}
so we can obtain
\begin{equation*}
	\begin{aligned}
		E_r = \frac{u_s}{1+\xi'(s)(h(t)-h_\infty},\\
		I_r = \frac{v_s}{1+\xi'(s)(h(t)-h_\infty}.
	\end{aligned}
\end{equation*}
And if we differentiate both sides twice with respect to $s$,
\begin{equation*}
	\begin{aligned}
		&\xi''(s)(h(t)-h_\infty)E_r+(1+\xi'(s)(h(t)-h_\infty)\frac{\partial}{\partial s}E_r=u_{ss},\\
		&\xi''(s)(h(t)-h_\infty)E_r+(1+\xi'(s)(h(t)-h_\infty)\frac{\partial r}{\partial s}\frac{\partial}{\partial r}E_r=u_{ss},\\
		&\xi'(s)(h(t)-h_\infty)E_r+(1+\xi'(s)(h(t)-h_\infty))^2E_{rr}=u_{ss}.
	\end{aligned}
\end{equation*}
Similarly,
\begin{equation*}
	\xi''(s)((h(t)-h_\infty)I_r+(+\xi'(s)(h(t)-h_\infty))^2I_{rr}=v_{ss},
\end{equation*}
so we can obtain
\begin{equation*}
	\begin{aligned}
		E_{rr}=\frac{u_{ss}-\xi''(s)(h(t)-h_\infty)E_r}{(1+\xi'(s)(h(t)-h_\infty))^2}=\frac{(1+\xi'(s)(h(t)-h_\infty))u_{ss}-\xi''(s)(h(t)-h_\infty)u_s}{(1+\xi'(s)(h(t)-h_\infty))^3},
	\end{aligned}
\end{equation*}
and 
\begin{equation*}
	I_{rr}= \frac{v_{ss}-\xi''(s)(h(t)-h_\infty)I_r}{(1+\xi'(s)(h(t)-h_\infty))^2}=\frac{((1+\xi'(s))(h(t)-h_\infty))v_{ss}-\xi''(s)(h(t)-h_\infty)u_s}{(1+\xi'(s)(h(t)-h_\infty))^3}.
\end{equation*}
In other case, if we differentiate both sides with respect to $t$,
\begin{equation*}
	\begin{aligned}
		&\xi(s)h'(t)E_r+E_t=u_t,\\
		&\frac{\xi(s)h'(t)}{1+\xi'(s)(h(t)-h_\infty)}u_s+E_t=u_t,
	\end{aligned}
\end{equation*}
and
\begin{equation*}
	\frac{\xi(s)h'(t)}{1+\xi'(s)(h(t)-h_\infty)}v_s+I_t=v_t,
\end{equation*}
so we can obtain
\begin{equation*}
	\begin{aligned}
		E_t=u_t-\frac{\xi(s)h'(t)}{1+\xi'(s)(h(t)-h_\infty)}u_s,\\
		I_t=v_t-\frac{\xi(s)h'(t)}{1+\xi'(s)(h(t)-h_\infty)}v_s.
	\end{aligned}
\end{equation*}
By using these calculations, we can change of the eqution\cref{free boundary SEIS} of the Exposed $E$,
\begin{equation}\label{modified free boundary SEIS3}
	\begin{aligned}
		E_t-d\Delta_rE&=u_t-\frac{\xi(s)h'(t)}{1+\xi'(s)(h(t)-h_\infty)}u_s\\
		&\quad -d\left(\frac{(1+\xi'(s)(h(t)-h_\infty))u_{ss}-\xi''(s)(h(t)-h_\infty)u_s}{(1+\xi'(s)(h(t)-h_\infty))^3}\right.\\
		&\quad \left.+\frac{n-1}{r}\frac{u_s}{1+\xi'(s)(h(t)-h_\infty)}\right)\\
		&=u_t-\frac{\xi(s)h'(t)}{1+\xi'(s)(h(t)-h_\infty)}u_s-\frac{du_{ss}}{(1+\xi'(s)(h(t)-h_\infty))^2}\\
		&\quad +\frac{d\xi''(s)(h(t)-h_\infty)u_s}{(x+\xi'(s)(h(t)-h_\infty))^3}-d\frac{n-1}{r}\frac{u_s}{1+\xi'(s)(h(t)-h_\infty)}\\
		&=u_t-\frac{\xi(s)h'(t)}{1+\xi'(s)(h(t)-h_\infty)}u_s-\frac{d\Delta_s u}{(1+\xi'(s)(h(t)-h_\infty))^2}\\
		&\quad +\frac{d\xi''(s)(h(t)-h_\infty)u_s}{(1+\xi'(s)(h(t)-h_\infty))^3}\\
		&\quad +\frac{d(n-1)u_s(-s\xi'(s)+\xi(s))(h(t)-h_\infty)}{s(s+\xi(s)(h(t)-h_\infty)(1+\xi'(s)(h(t)-h_\infty))^2}.
	\end{aligned}
\end{equation}
If we define
\begin{equation*}
	\begin{aligned}
		&\sqrt{X(h(t),s)}\coloneqq \frac{\partial s}{\partial r}=\frac{1}{1+\xi'(s)(h(t)-h_0)},
			\\
		&Y(h(t),s)\coloneqq \frac{\partial^2s}{\partial r^2}=-\frac{\xi''(s)(h(t)-h_0)}{[1+\xi'(s)(h(t)-h_0)]^3},
			\\
		&Z(h(t),s)\coloneqq -\frac{1}{h'(t)}\frac{\partial s}{\partial t}=\frac{\xi(s)}{1+\xi'(s)(h(t)-h_0)},
			\\
		&W(h(t),s)\coloneqq(n-1)X\frac{(s\xi'-\xi)(h(t)-h_0)}{s(s+\xi(s)(h(t)-h_0))}.\,
	\end{aligned}
\end{equation*}
then we can rewrite the equation\cref{modified free boundary SEIS3}
\begin{equation*}
	u_t-Xd\Delta_s u -(Yd+Zh'(t)+Wd)u_s=(1-p)\alpha vw-(\beta_1+r_1+\mu_2)u.
\end{equation*}
Similarly, we had a direct calculation for $I$ and it is changed to
\begin{equation*}
	\begin{aligned}
		I_t-d_3\Delta_r I=&v_t-\frac{\xi(s)h'(t)}{1+\xi'(s)(h(t)-h_\infty)}v_s\\
		&-d\left( \frac{(1+\xi'(s)(h(t)-h_\infty))v_{ss}-\xi''(s)(h(t)-h_\infty)v_s}{(1+\xi'(s)(h(t)-h_\infty))^3} \right. \\
		&\left. \frac{n-1}{r}\frac{v_s}{1+\xi'(s)(h(t)-h_\infty)} \right),
	\end{aligned}
\end{equation*}
so
\begin{equation*}
	I_t-d\Delta_r I = v_t-Xd\Delta_s v-(Yd+Zh'(t)+Wd)v_s=p\alpha vw+\beta_1u-(r_2+\mu_3)v.
\end{equation*}
Therefore, the equations become
\begin{equation*}
	\left \{ \begin{aligned}
		&u_t-Xd\Delta_s u-(Yd+Zh'(t)+Wd)u_s=(1-p)\alpha vw-(\beta_1+r_1+\mu_2)u,\\
		&v_t-Xd\Delta_s v-(Yd+Zh'(t)+Wd)v_s=p\alpha vw+\beta_1 u-(r_2+\mu_3)v,
	\end{aligned} \right.
\end{equation*}
for $(s,t)\in \Omega_{h_\infty}$, and
\begin{equation*}
	\left \{ \begin{aligned}
		&h'(t)=-\beta v_s(h_0,t)-\mu w_s(h_0,t), \quad &&s=h_\infty,\\
		&u(s,t)=v(s,t)=0, &&s\geq h_\infty.
	\end{aligned} \right.
\end{equation*}
\end{proof}

\end{document}
